\documentclass[12pt]{article}

\usepackage[margin=1in]{geometry}
\usepackage{amsmath,amsthm,amssymb,scrextend}
\usepackage{fancyhdr}
\usepackage{graphicx}
\usepackage{enumitem}
\pagestyle{fancy}
\setlength{\headheight}{28pt}

% Red "alertblock"-style boxes (Beamer-like)
\usepackage{xcolor}
\usepackage[most]{tcolorbox}
\newtcolorbox{alertblock}[1]{%
  colback=red!5!white,
  colframe=red!70!black,
  fonttitle=\bfseries,
  title=#1,
  sharp corners,
  boxrule=0.9pt,
  left=1.2mm,right=1.2mm,top=1mm,bottom=1mm
}

% Common shortcuts (kept close to template)
\newcommand{\R}{\mathbb R}
\newcommand{\N}{\mathbb N}
\newcommand{\Z}{\mathbb Z}
\newcommand{\E}{\mathbb E}
\newcommand{\Var}{\mathrm{Var}}
\newcommand{\imp}{\Rightarrow}
\newcommand{\set}[1]{\left\{ #1 \right\}}
\DeclareMathOperator*{\plim}{plim}

\usepackage{datetime}

% Date (update as needed)
\newdate{date}{27}{03}{2026}

\usepackage[round]{natbib}
\usepackage[colorlinks=true, pdfstartview=FitV, linkcolor=blue,
citecolor=blue, urlcolor=blue]{hyperref}

\newtheorem*{definition}{Definition}
\theoremstyle{definition}
\newtheorem{problem}{Problem}
\theoremstyle{plain}
\newtheorem{exercise}{Exercise}

\begin{document}

\lhead{TA: Rafael Lincoln \\ Prof. Tim Cogley}
\rhead{Econometrics II (Time Series) \\ For: \displaydate{date}}

\section*{Problem Set 5 --- GMM: Identification, Testing, Misspecification, and Optimal Instruments}

\begin{alertblock}{Important}
We \textbf{do not collect} this problem set, and you \textbf{do not need to do all exercises}.
We will cover the solutions in class and provide written solutions afterwards for your study.
The goal is to give you more material to prepare for the exam.
\end{alertblock}

\vspace{0.8em}

% ============================================================================
% PROBLEM 1: 2018 Practice Final - EHS (1988) testing subset of moments
% ============================================================================
\begin{problem}[Testing a subset of moment conditions: Eichenbaum--Hansen--Singleton (1988)]
\label{prob:ehs}

\citet{EichenbaumHansenSingleton1988} estimate a representative-agent model of consumption and leisure choice. The agent chooses consumption and leisure streams $(c_t, \ell_t)_{t \geq 0}$ to maximize the time- and state-separable utility function
\[
\E\left[ \sum_{t=0}^{\infty} \beta^t u(c_t, \ell_t; \gamma, \alpha) \right]
\]
subject to a budget constraint. This leads to first-order conditions that can be written as
\[
\E\left[ h_1\!\left( \frac{c_{t+1}}{c_t}, \frac{\ell_t}{c_t}; \beta, \gamma, \alpha \right) R_{t+1} \mid \mathcal{F}_t \right] = 1, \qquad
\E\left[ h_2\!\left( \frac{c_{t+1}}{c_t}, \frac{\ell_t}{c_t}; \beta, \gamma, \alpha \right) R_{t+1} \mid \mathcal{F}_t \right] = 1,
\]
where $h_1$ and $h_2$ are known functions, $(\beta, \gamma, \alpha)$ are unknown preference parameters in some parameter space $\Theta \subset \R^3$, and $R_{t+1}$ is the return on some asset. The authors estimate the true preference parameters by GMM using the population moment conditions
\[
\E\left[ \begin{pmatrix} h_1(\cdot; \beta, \gamma, \alpha) R_{t+1} - 1 \\ h_2(\cdot; \beta, \gamma, \alpha) R_{t+1} - 1 \end{pmatrix} \otimes z_t \right] = 0,
\]
where $z_t$ is a $4 \times 1$ vector of variables in $\mathcal{F}_t$. You may assume that the data $w_1, \ldots, w_n$ are strictly stationary and ergodic (with $w_t$ containing $c_{t+1}/c_t$, $\ell_t/c_t$, $R_{t+1}$, and $z_t$), and that all usual regularity conditions hold.

The presence of labor market frictions may make the moment conditions involving $h_2$ invalid, i.e.\ it may be that
\[
\E\left[ \bigl( h_2(\cdot; \beta_0, \gamma_0, \alpha_0) R_{t+1} - 1 \bigr) \otimes z_t \right] \neq 0.
\]

Describe how you would test whether the moment conditions involving $h_2$ are invalid. In your answer, state clearly: (i) the null and alternative hypotheses, (ii) the test statistic you would use, and (iii) its asymptotic distribution under the null.
\end{problem}

\vspace{0.8em}

% ============================================================================
% PROBLEM 2: ps4.pdf Q1 - Testing with GMM, omit 1/2 factor
% ============================================================================
\begin{problem}[GMM objective and test statistics]
\label{prob:gmm-half}

Suppose you omit the $\frac{1}{2}$ factor in the usual definition of the GMM objective function and instead use
\[
\tilde{Q}_n(\theta) = -\bar{g}_n(\theta)' \hat{W} \bar{g}_n(\theta),
\]
where $\bar{g}_n(\theta) = \frac{1}{n}\sum_{t=1}^n g(x_t; \theta)$. Suppose you have an optimal weighting matrix (i.e.\ $\hat{W} \xrightarrow{p} S^{-1}$).

\begin{enumerate}[label=(\alph*)]
\item State the asymptotic distribution of the Wald, LM, and QLR tests of $H_0\colon a(\theta_0) = 0$ that you would obtain using this modified objective $\tilde{Q}_n$.
\item Explain how the conclusions of your tests could be affected by omitting the $\frac{1}{2}$ factor. You may assume that all regularity conditions hold.
\end{enumerate}

\noindent\textit{Hint:} This need not be long; inspect how the test statistics depend on the objective and how omitting the factor affects their asymptotic distribution.
\end{problem}

\vspace{0.8em}

% ============================================================================
% PROBLEM 3: ps3.pdf Q3 - Misspecified GMM
% ============================================================================
\begin{problem}[Misspecified GMM]
\label{prob:misspec-gmm}

Suppose you are estimating an over-identified model by GMM, but the model is misspecified: there is no $\theta \in \Theta$ such that $\E[g(x_t; \theta)] = 0$.

\begin{enumerate}[label=(\alph*)]
\item What is the pseudo-true parameter that you are in fact estimating? Show how it depends on your choice of weighting matrix.
\item You may assume the pseudo-true parameter is unique and that your estimator is consistent for it. Mimic the proof of asymptotic normality for GMM. Where does the argument break down?
\item Now suppose the model is ``locally'' misspecified: the degree of misspecification is of order sampling error,
\[
\E[g(x_t; \theta_0)] = \frac{c}{\sqrt{n}},
\]
where $\theta_0$ is the pseudo-true parameter and $c$ is a vector of constants. Derive the asymptotic distribution of your estimator. How does it differ from the correctly specified case?
\item Are the usual asymptotic variance estimators valid in this case? Explain.
\end{enumerate}

\noindent\textit{Remark:} Similar local asymptotic frameworks are used when assumptions only approximately hold (e.g.\ ``weak identification''). The asymptotic distribution of $\hat{\theta}$ under global misspecification is more complicated and depends on how the weighting matrix and gradient are estimated.
\end{problem}

\vspace{0.8em}

% ============================================================================
% PROBLEM 4: Lista de exercícios Q10 - GMM (translated)
% ============================================================================
\begin{problem}[GMM: identification, efficient weighting, and optimal instruments]
\label{prob:gmm-optimal}

Consider estimation of a parameter vector $\beta \in \R^k$ by the Generalized Method of Moments (GMM) applied to the conditional moment condition $\E[e(z_i; \beta_0) \mid x_i] = 0$, where $e(z_i; \beta) \in \R$ is continuous and differentiable in $\beta$. Suppose $u_i = (x_i', z_i')'$ is a vector of i.i.d.\ random variables.

The conditional moment above implies the unconditional moment $\E[e(z_i; \beta_0) A(x_i)] = 0$ for any vector function $A(x_i) \in \R^m$ of $x_i$. Define $g(u_i; \beta) = e(z_i; \beta) A(x_i)$. The GMM estimator is
\[
\hat{\beta} = \arg\min_{\beta} \; \bar{g}(\beta)' W_n \bar{g}(\beta), \qquad \bar{g}(\beta) = \frac{1}{n} \sum_{i=1}^n g(u_i; \beta),
\]
where $W_n$ is an $(m \times m)$ symmetric positive definite matrix. Under standard regularity conditions, $\sqrt{n}(\hat{\beta} - \beta_0) \xrightarrow{d} N(0, \Sigma)$, where the asymptotic covariance matrix is
\[
\Sigma = (G' W G)^{-1} G' W V W G (G' W G)^{-1},
\]
with $W = \plim W_n$, $G = \E\bigl[ \frac{\partial}{\partial \beta} g(u_i; \beta_0) \bigr]$ $(m \times k)$, and $V = \E\bigl[ A(x_i) e(z_i; \beta_0)^2 A(x_i)' \bigr]$ $(m \times m)$.

\begin{enumerate}[label=(\alph*)]
\item What is a sufficient condition for local identification of $\beta_0$?
\item Find the \emph{infeasible} weighting matrix $W_n$ (the one that depends on the true $\beta_0$) that yields efficient GMM. Justify asymptotic optimality: obtain the asymptotic covariance matrix $\Sigma^\#$ for $\sqrt{n}(\hat{\beta} - \beta_0)$ under this choice and show that $\Sigma - \Sigma^\# \geq 0$ (i.e.\ $\Sigma - \Sigma^\#$ is positive semidefinite) for any other choice of $W$.
\item The GMM estimator depends on the choice of $A(x_i)$. Find the asymptotic covariance $\Sigma^*$ when
\[
A(x_i)^* = \frac{D(x_i)}{\sigma^2(x_i)}, \quad \text{where} \quad D(x_i) = \E\!\left[ \frac{\partial e(z_i; \beta)}{\partial \beta} \,\bigg|\, x_i \right]_{\beta = \beta_0}, \quad \sigma^2(x_i) = \E\bigl[ e(z_i; \beta_0)^2 \mid x_i \bigr].
\]
\item Show that $\Sigma - \Sigma^* \geq 0$ (with strict inequality in the matrix sense when the optimal $A^*$ is not used).
\item How would you compute $A^*(x_i)$ in practice?
\end{enumerate}
\end{problem}

\vspace{0.8em}

% ----------------------------------------------------------------------------
\begin{thebibliography}{1}
\bibitem[Eichenbaum et al.(1988)]{EichenbaumHansenSingleton1988}
Eichenbaum, M., L.\ P.\ Hansen, and K.\ J.\ Singleton (1988).
\newblock ``A time series analysis of representative agent models of consumption and leisure choice under uncertainty.''
\newblock \emph{Quarterly Journal of Economics} 103(1), 51--78.
\end{thebibliography}

\end{document}
